% !TeX root = ../../Book3.tex
\chapter{Dirichlet 问题与 Perron 方法}
\label{b3:ch:16}
\BookChapterNumber{b3:ch:16}
% 本章摘要（不计入编号）
\begin{chapterabstract}
圆盘 Poisson 积分可以直接从边值计算内部函数。一般区域的解则由 Perron 次函数族构造，是否达到每个边界值取决于障碍函数。本章完成包络的调和性证明，并在正则域上把解表示为边界调和测度的积分。
\end{chapterabstract}

% 本章目标（不计入编号）
\begin{learninggoals}
能计算圆盤 Poisson 延拓；能说明 Perron 族的调和提升为何保留次函数条件；能构造外锥障碍并检验边界极限；能在给定边界延拓条件下使用调和测度的共形不变性。
\end{learninggoals}

在单位圆周给定 \(\cos(3t)\)，内部解是 \(r^3\cos(3\theta)\)。如果把圆盘中心挖去，再独立给这个孤立边界点赋值，解却未必接受该数据。一个边值问题因而包含两项责任：构造内部调和函数，证明它达到边界。前者由平均值、比较原理和 Harnack 收敛解决，后者由障碍解决。

\ProseParagraphBreak

本章使用\chapref{b3:ch:15}的次调和比较与 Harnack 原理。圆盘核的计算可与 Axler、Bourdon、Ramey 的《Harmonic Function Theory》\cite{Axler2001Harmonic}的本地背景材料比较；一般域的势论继续阅读 Ransford 的《Potential Theory in the Complex Plane》\cite{Ransford1995Potential}。

% !TeX root = ../../Book3.tex
\section{Dirichlet 问题}
\label{b3:sec:1601}

给定边界温度，怎样确定内部平衡状态？第十五章已经用 Poisson 积分解决了圆盘上的连续边值问题。本节从这个解出发，研究边界误差怎样传入内部、不同频率怎样衰减，以及边值不连续时应当采用什么极限。上半平面的公式还会说明，无界区域的唯一性需要另加增长限制。

\subsection{边值问题}
\label{b3:subsec:160101}

设 \(D\subseteq\mathbb C\) 为有界区域，\(\varphi\in C(\partial D;\mathbb R)\)。寻找一个 \(u\in C(\overline D;\mathbb R)\)，使
\[
 \Delta u=0\quad\text{于 }D,
 \quad
 u=\varphi\quad\text{于 }\partial D,
\]
称为给定数据 \(\varphi\) 的\term[dirichlet-wenti]{Dirichlet 问题}[Dirichlet problem]。这里方程要求 \(u|_D\in C^2(D)\)，边界条件则按闭包上的连续性理解；满足这些要求的函数称为经典解。域内没有热源对应 Laplace 方程，边界温度则选择众多调和函数中的一个。

\ProseParagraphBreak
唯一性与存在性承担不同的任务。若 \(u,v\) 都是同一数据的经典解，则 \(u-v\) 在有界域 \(D\) 调和，连续到 \(\overline D\)，边界值为零。对差及其负号应用最大值原理，得到 \(u=v\)。这个论证不要求边界光滑，却没有构造任何解。去心圆盘上在孤立边界点另指定一个温度，就可能得不到经典解；下一节将给出具体证明。

\ProseParagraphBreak
比较原理同时给出数据扰动的控制。设 \(\varphi,\psi\in C(\partial D;\mathbb R)\) 分别已有经典解 \(u,v\)，并令
\[
 d=\sup_{\partial D}|\varphi-\psi|.
\]
函数 \(u-v-d\) 与 \(v-u-d\) 在边界非正，因此
\[
 \sup_D|u-v|\le d.
\]
若 \(\varphi\le\psi\)，对 \(u-v\) 直接比较还得 \(u\le v\)。所以一旦存在性成立，边界数据的最大误差不会在内部放大。特别地，若一列可解的连续边值一致收敛，则其解在整个闭包上一致 Cauchy；极限的内部调和性可由圆周平均等式传给极限，再用\cref{b3:thm:15-mean-value}的连续函数逆命题得到。

\ProseParagraphBreak
边界上的同一个公式，不能不加检验地用作内部函数。例如单位圆周上的 \(|z|^2\) 恒等于一，故其调和延拓是常数一；函数 \(|z|^2=x^2+y^2\) 本身满足 \(\Delta|z|^2=4\)，并不是这份边值的解。边界条件只规定圆周上的值，域内公式还必须满足方程。

\ProseParagraphBreak
实值连续数据是经典问题的基本范围。复值数据可以分别延拓实部、虚部；有界但不连续的数据仍可作 Poisson 积分，不过此时必须逐点说明边界极限，不能预先要求解属于 \(C(\overline D)\)。

\subsection{圆盘 Poisson 核}
\label{b3:subsec:160102}

圆盘中的构造要同时满足两件事：对内部点保持调和性，对接近的边界点集中权重。设 \(R>0\)，\(|z|<R\)。第十五章引入的 Poisson 核为
\[
 P_R(z,t)=\frac{R^2-|z|^2}{|Re^{\mathrm i t}-z|^2}
 =\operatorname{Re}\frac{Re^{\mathrm i t}+z}{Re^{\mathrm i t}-z}.
\]
右端是关于 \(z\) 的全纯函数的实部，因而调和；左端直接显示它为正。分子中的 \(R^2-|z|^2\) 还给出正确的归一化：
\[
 P_R(z,t)=1+2\operatorname{Re}\sum_{n\ge1}
       \left(\frac zR\right)^n e^{-\mathrm i nt},
 \quad
 \frac1{2\pi}\int_{-\pi}^{\pi}P_R(z,t)\,\mathrm dt=1.
\]
几何级数在每个较小闭盘上关于 \(t\) 一致收敛，所以可逐项积分；所有非零频率的积分均为零。保持常数与集中质量是核逼近边值时各自需要的性质。

\ProseParagraphBreak
把 \(z=re^{\mathrm i\theta}\)、\(\rho=r/R\) 代入，可把半径尺度从角变量中分离出来。记
\[
 p_\rho(s)\coloneqq
 \frac{1-\rho^2}{1-2\rho\cos s+\rho^2},
 \quad 0\le\rho<1.
\]
则 \(P_R(re^{\mathrm i\theta},t)=p_\rho(t-\theta)\)，而 \(p_\rho\) 为偶的 \(2\pi\) 周期函数。接近边界时，控制它的尺度是 \(\varepsilon=1-\rho\)，即到边界的径向距离除以 \(R\)。

\Needspace{6\baselineskip}
\begin{lemma}[Poisson 核的集中估计]\label{b3:lem:16-poisson-localization}
设 \(1/2\le\rho<1\)，\(\varepsilon=1-\rho\)，并令
\(p_\rho(t)=(1-\rho^2)/(1-2\rho\cos t+\rho^2)\)。对 \(|t|\le\pi\)，有
\[
 \frac{\varepsilon}{\varepsilon^2+t^2}
 \le p_\rho(t)
 \le\frac{\pi^2\varepsilon}{\varepsilon^2+t^2}.
\]
特别地，对 \(0<\delta\le\pi\)，
\[
 \frac1{2\pi}\int_{\delta\le|t|\le\pi}p_\rho(t)\,\mathrm dt
 \le\frac{\pi\varepsilon}{\delta}.
\]
\end{lemma}
\begin{proof}
分母可写为
\[
 1-2\rho\cos t+\rho^2
 =\varepsilon^2+4\rho\sin^2(t/2).
\]
在 \(|t|\le\pi\) 上，\(|t|/\pi\le|\sin(t/2)|\le|t|/2\)，故
\[
 \frac2{\pi^2}(\varepsilon^2+t^2)
 \le\varepsilon^2+4\rho\sin^2(t/2)
 \le\varepsilon^2+t^2.
\]
分子为 \(\varepsilon(1+\rho)\)，介于 \(\varepsilon\) 与 \(2\varepsilon\) 之间，得到双边估计。利用偶性和 \(t\ge\delta\) 时 \(\varepsilon^2+t^2\ge t^2\)，进一步有
\[
 \frac1{2\pi}\int_{\delta\le|t|\le\pi}p_\rho(t)\,\mathrm dt
 \le\pi\varepsilon\int_\delta^\infty\frac{\mathrm dt}{t^2}
 =\frac{\pi\varepsilon}{\delta}.
\]
\end{proof}

在中心 \(\rho=0\) 时，核恒等于一。接近边界时，它在 \(t=0\) 的高度为 \((1+\rho)/(1-\rho)\)，在 \(|t|\) 与 \(1-\rho\) 同阶的角区间中仍有相同量级。高度增大、宽度缩小，而总积分保持 \(2\pi\)；任何固定远弧的质量则由上面的估计趋于零。这个尺度关系将同时给出边界收敛速度和跳跃点的极限。

\subsection{Poisson 积分}
\label{b3:subsec:160103}

用正的归一化核对边界数据加权，便得到圆盘解。下面保留一个可直接引用的 Dirichlet 问题版本；圆盘上的存在性与边界连续性的证明已经在上一章完成。

\Needspace{6\baselineskip}
\begin{theorem}[圆盘 Dirichlet 问题]\label{b3:thm:16-poisson}
设 \(R>0\)，\(\varphi:\partial B(0,R)\to\mathbb R\) 连续，并令
\[
 P_R(z,t)=\frac{R^2-|z|^2}{|Re^{\mathrm i t}-z|^2},
 \quad |z|<R.
\]
则存在唯一 \(u\in C(\overline{B(0,R)};\mathbb R)\)，在 \(B(0,R)\) 调和且在边界等于 \(\varphi\)。它在内部由
\[
 u(z)=\frac1{2\pi}\int_0^{2\pi}
 P_R(z,t)\varphi(Re^{\mathrm i t})\,\mathrm dt
\]
给出。
\end{theorem}
\begin{proof}
\cref{b3:thm:15-poisson-continuous}适用于任意正半径圆盘上的实值连续圆周数据，正是这里的条件。该定理用内部紧集上的核导数界证明积分调和，用近弧上的连续性和远弧上的质量消失证明任意方向的边界连续性，再用有界域最大值原理证明唯一性。因此把其 Poisson 积分用于当前 \(\varphi\)，便得到所述经典解。
\end{proof}

记这一延拓为 \(\mathcal P_R\varphi\)。积分直接显示它对数据线性、保持常数并保持大小关系，且
\[
 |(\mathcal P_R\varphi)(z)|
 \le\|\varphi\|_\infty,
 \quad
 (\mathcal P_R\varphi)(0)
 =\frac1{2\pi}\int_0^{2\pi}\varphi(Re^{\mathrm i t})\,\mathrm dt.
\]
中心只看到边界平均，其余角向信息随着离开中心逐渐出现。为精确描述这种变化，先把边界数据与各个角频率配对。

\Needspace{6\baselineskip}
\begin{proposition}[角频率的延拓与衰减]\label{b3:prop:16-poisson-frequencies}
设 \(R>0\)，\(\varphi:\partial B(0,R)\to\mathbb C\) 为有界 Borel 可测函数，\(u\) 由其 Poisson 积分定义。令
\[
 c_n=\frac1{2\pi}\int_0^{2\pi}
       \varphi(Re^{\mathrm i t})e^{-\mathrm i nt}\,\mathrm dt,
 \quad n\in\mathbb Z.
\]
则对 \(0\le r<R\)，
\[
 u(re^{\mathrm i\theta})
 =\sum_{n\in\mathbb Z}c_n
       \left(\frac rR\right)^{|n|}e^{\mathrm i n\theta}.
\]
该级数在每个内部闭盘上绝对一致收敛，且可按实变量 \(x,y\) 逐项求任意阶导数。其和的实部与虚部均调和。若 \(0<q<1\)、\(N\ge0\) 为整数，则截断到 \(|n|\le N\) 后，在 \(|z|\le qR\) 上的误差不超过
\[
 \frac{2\|\varphi\|_\infty q^{N+1}}{1-q}.
\]
\end{proposition}
\begin{proof}
有界性给出 \(|c_n|\le\|\varphi\|_\infty\)。核的几何级数展开在 \(|z|\le qR\) 上关于 \(t\) 绝对一致收敛，乘以有界可测的边界数据后仍可交换求和与积分。于是
\[
 u(z)=c_0+\sum_{n\ge1}
 \left[c_n\left(\frac zR\right)^n+
       c_{-n}\left(\frac{\overline z}R\right)^n\right].
\]
这与角变量表达相同。在 \(|z|\le qR\) 上，取 \(k\) 阶实偏导数后，每个 \(n\ge k\) 的项可由常数乘以 \(n^kq^{n-k}\) 控制；这个数列可求和。有限个低次项也可直接求导。因此各阶导数级数一致收敛，每项的实部、虚部都调和，故和也调和。最后对 \(|n|>N\) 的项取绝对值并求几何级数，得到所述余项界。
\end{proof}

系数 \(c_n\) 称为边界数据的\term[fourier-xishu]{Fourier 系数}[Fourier coefficient]。上面的推导只用核的几何级数，不要求边界函数自身的 Fourier 级数在任何边界点收敛。因子 \((r/R)^{|n|}\) 已使内部展开绝对收敛；当 \(r/R\) 固定小于一时，较高频率受到更强的压制。

\ProseParagraphBreak
对 \(n\ge1\)，数据 \(\cos(nt)\)、\(\sin(nt)\) 分别延拓为
\[
 \left(\frac rR\right)^n\cos(n\theta),
 \quad
 \left(\frac rR\right)^n\sin(n\theta).
\]
正频率 \(e^{\mathrm i nt}\) 对应 \((z/R)^n\)，负频率对应 \((\overline z/R)^{|n|}\)。实数据满足 \(c_{-n}=\overline{c_n}\)，两部分配对成为实函数；调和延拓并不要求负频率消失。

\ProseParagraphBreak
例如圆周上的数据 \(\varphi(z)=(\operatorname{Re}z)^2\) 可写成
\[
 \varphi(Re^{\mathrm i t})
 =R^2\cos^2t=\frac{R^2}{2}(1+\cos2t).
\]
它的调和延拓因而为
\[
 u(x+\mathrm i y)=\frac{R^2+x^2-y^2}{2}.
\]
在圆周上由 \(x^2+y^2=R^2\) 得 \(u=x^2\)，内部则有 \(\Delta u=0\)。中心值 \(R^2/2\) 是边界平均，不是把边界公式 \(x^2\) 直接代入中心所得的零。

\subsection{Poisson 积分的边界行为与调和测度}
\label{b3:subsec:160104}

连续边值已经保证任意方向的边界连续性。沿同心圆趋近时，核的尺度估计还能把数据的连续模转成统一误差，说明光滑程度怎样影响逼近速度。

\Needspace{6\baselineskip}
\begin{proposition}[径向边界误差]\label{b3:prop:16-poisson-boundary-rate}
设 \(R>0\)，\(\varphi:\partial B(0,R)\to\mathbb R\) 连续，\(u=\mathcal P_R\varphi\)。定义角变量的连续模
\[
 \omega_\varphi(\delta)=
 \sup_{\substack{\theta\in\mathbb R\\|t|\le\delta}}
 |\varphi(Re^{\mathrm i(\theta+t)})
       -\varphi(Re^{\mathrm i\theta})|,
 \quad 0<\delta\le\pi.
\]
若 \(R/2\le r<R\)，\(\varepsilon=1-r/R\)，则有绝对常数 \(C\)，使
\[
 \sup_\theta|u(re^{\mathrm i\theta})-\varphi(Re^{\mathrm i\theta})|
 \le\omega_\varphi(\delta)
       +C\|\varphi\|_\infty\frac{\varepsilon}{\delta}.
\]
因此左端随 \(r\uparrow R\) 趋于零。若进一步有
\[
 |\varphi(Re^{\mathrm i(\theta+t)})
       -\varphi(Re^{\mathrm i\theta})|
 \le H|t|^\alpha
 \quad(\theta\in\mathbb R,\ |t|\le\pi),
\]
其中 \(H\ge0\)，则当 \(0<\alpha<1\) 时，上述左端至多为
\(C_\alpha H\varepsilon^\alpha\)；当 \(\alpha=1\) 时，至多为
\(CH\varepsilon\log(e/\varepsilon)\)。常数与 \(R,r,\varphi,H\) 无关，\(C_\alpha\) 可依赖 \(\alpha\)。
\end{proposition}
\begin{proof}
写 \(\rho=r/R\)。核的归一化与周期性给出
\[
 u(re^{\mathrm i\theta})-\varphi(Re^{\mathrm i\theta})
 =\frac1{2\pi}\int_{-\pi}^{\pi}p_\rho(t)
 \bigl[\varphi(Re^{\mathrm i(\theta+t)})
       -\varphi(Re^{\mathrm i\theta})\bigr]\,\mathrm dt.
\]
在 \(|t|<\delta\) 上，边值差不超过 \(\omega_\varphi(\delta)\)，核质量至多一；其余部分用边值差至多 \(2\|\varphi\|_\infty\) 及\cref{b3:lem:16-poisson-localization}的远弧估计，得到第一式。连续函数在紧圆周上一致连续，故先使 \(\omega_\varphi(\delta)\) 小，再固定 \(\delta\) 令 \(\varepsilon\downarrow0\)，即得一致收敛。

有 Hölder 条件时，不再先固定一条近弧，而在整个积分中保留差值大小。核的上界给
\[
 |u(re^{\mathrm i\theta})-\varphi(Re^{\mathrm i\theta})|
 \le CH\int_0^\pi
        \frac{\varepsilon t^\alpha}{\varepsilon^2+t^2}\,\mathrm dt.
\]
若 \(0<\alpha<1\)，代入 \(t=\varepsilon s\)，右侧至多为
\[
 CH\varepsilon^\alpha
       \int_0^\infty\frac{s^\alpha}{1+s^2}\,\mathrm ds.
\]
该积分在零附近由 \(s^\alpha\) 控制，在无穷远由 \(s^{\alpha-2}\) 控制，两端均可积。若 \(\alpha=1\)，直接计算得到
\[
 \int_0^\pi\frac{\varepsilon t}{\varepsilon^2+t^2}\,\mathrm dt
 =\frac{\varepsilon}{2}
     \log\left(1+\frac{\pi^2}{\varepsilon^2}\right)
 \le C\varepsilon\log(e/\varepsilon),
\]
于是得到端点估计。
\end{proof}

Lipschitz 情形的对数因子一般不能删去。取单位圆周数据
\(\varphi(e^{\mathrm i t})=|t|\)，其中 \(-\pi\le t\le\pi\)，并作周期延拓。它是到角度零的圆周距离，满足角变量的 Lipschitz 条件。在点 \(1\) 的径向线上，令 \(r=1-\varepsilon\)、\(0<\varepsilon\le1/2\)，则
\[
 (\mathcal P_1\varphi)(r)-\varphi(1)
 =\frac1\pi\int_0^\pi p_r(t)t\,\mathrm dt
 \ge\frac{\varepsilon}{2\pi}
       \int_\varepsilon^1\frac{\mathrm dt}{t}
 =\frac{\varepsilon}{2\pi}\log(1/\varepsilon).
\]
这里使用了核的下界及 \(t\ge\varepsilon\) 时
\(\varepsilon^2+t^2\le2t^2\)。因此对所有 Lipschitz 边值声称误差为 \(O(1-r)\) 并不成立。

\ProseParagraphBreak
固定 \(R>0\) 与 \(z\in B(0,R)\)，把核当作边界上的密度，定义概率测度
\[
 \mathrm d\omega^z(Re^{\mathrm i t})
 =\frac1{2\pi}P_R(z,t)\,\mathrm dt.
\]
它称为圆盘的\term[diaohe-cedu]{调和测度}[harmonic measure]，满足
\[
 (\mathcal P_R\varphi)(z)=\int_{\partial B(0,R)}\varphi\,\mathrm d\omega^z.
\]
中心处它等于归一化弧长；非中心处则对较近的边界赋予较大权重。这里的测度由显式密度定义，稍后的\cref{b3:sec:1604}将在一般正则域上从解算子构造它。

\ProseParagraphBreak
对有界 Borel 数据，同一积分仍有意义，并由\cref{b3:prop:16-poisson-frequencies}给出内部调和函数。改变有限个边界点的取值不会改变积分，因为每个点的弧长测度为零。在任何一个数据连续的边界点 \(\xi\)，积分仍任意方向趋于 \(\varphi(\xi)\)：近 \(\xi\) 的弧上只需该点的连续性，远弧上的边值差由统一有界性控制，而核质量趋于零。全边界连续性的丢失，并不妨碍在其余连续点取到正确边值。

\ProseParagraphBreak
跳跃点处的情况取决于接近方向。下面直接用核的缩放计算这种差别，不需要预先知道某条圆弧的调和测度公式。

\Needspace{6\baselineskip}
\begin{proposition}[跳跃边值的方向极限]\label{b3:prop:16-poisson-jump}
设 \(\varphi:\partial B(0,1)\to\mathbb R\) 为有界 Borel 可测函数，\(u=\mathcal P_1\varphi\)。设在某个 \(\theta_0\in\mathbb R\) 处存在单侧极限
\[
 A_-=\lim_{t\uparrow0}\varphi(e^{\mathrm i(\theta_0+t)}),
 \quad
 A_+=\lim_{t\downarrow0}\varphi(e^{\mathrm i(\theta_0+t)}).
\]
对任意固定 \(\lambda\in\mathbb R\)，令
\(z_\varepsilon=(1-\varepsilon)e^{\mathrm i(\theta_0+\lambda\varepsilon)}\)，其中 \(0<\varepsilon<1/2\)。则
\[
 \lim_{\varepsilon\downarrow0}u(z_\varepsilon)
 =A_-\left(\frac12-\frac{\arctan\lambda}{\pi}\right)
  +A_+\left(\frac12+\frac{\arctan\lambda}{\pi}\right).
\]
特别地，径向极限为 \((A_-+A_+)/2\)。这些极限与
\(\varphi(e^{\mathrm i\theta_0})\) 本身的取值无关。
\end{proposition}
\begin{proof}
旋转后可设 \(\theta_0=0\)。取 \(0<\delta<\pi/2\)，把积分分成 \(|t|<\delta\) 与其余部分。当 \(\varepsilon\) 足够小时，核的中心角 \(\lambda\varepsilon\) 与零的角距离小于 \(\delta/2\)。远弧到中心的圆周角距离至少为 \(\delta/2\)，所以远弧贡献由集中估计趋于零。

在近弧作换元 \(t=\varepsilon s\)。对每个固定实数 \(s\)，有
\[
 \frac{\varepsilon}{2\pi}
 p_{1-\varepsilon}\bigl(\varepsilon(s-\lambda)\bigr)
 \longrightarrow\frac1{\pi\bigl(1+(s-\lambda)^2\bigr)}.
\]
这是把分子写成 \(\varepsilon(2-\varepsilon)\)，分母写成
\(\varepsilon^2+2(1-\varepsilon)(1-\cos(\varepsilon(s-\lambda)))\)
后，使用 \((1-\cos h)/h^2\to1/2\) 得到的。近弧中
\(|\varepsilon(s-\lambda)|<\pi\)，核上界又给出
\[
 \frac{\varepsilon}{2\pi}
 p_{1-\varepsilon}\bigl(\varepsilon(s-\lambda)\bigr)
 \le\frac{C}{1+(s-\lambda)^2}.
\]
将近弧积分的被积函数在 \(|s|\ge\delta/\varepsilon\) 处补零，右侧乘以 \(\|\varphi\|_\infty\) 就是整条实轴上的共同可积控制。对 \(s<0\)，数据趋于 \(A_-\)；对 \(s>0\)，趋于 \(A_+\)。支配收敛因此给出
\[
 \lim_{\varepsilon\downarrow0}u(z_\varepsilon)
 =\frac{A_-}{\pi}\int_{-\infty}^{0}
       \frac{\mathrm ds}{1+(s-\lambda)^2}
  +\frac{A_+}{\pi}\int_{0}^{\infty}
       \frac{\mathrm ds}{1+(s-\lambda)^2}.
\]
两项积分分别为 \(\pi/2-\arctan\lambda\) 与
\(\pi/2+\arctan\lambda\)，得到结论。单个角度的函数值不影响积分。
\end{proof}

当 \(A_-\ne A_+\) 时，改变 \(\lambda\) 就得到不同极限，因而解不能连续延拓到这个边界点。径向得到左右平均，是因为此时核关于该点对称；它并没有恢复一份连续边值。圆弧指示函数的端点属于这种情形，具体圆弧和实轴区间的权重将在\cref{b3:sec:1604}中计算。

\ProseParagraphBreak
虽然有限个点的指定数值不影响 Poisson 积分，但从“两个有界调和函数在其他边界点同值”推出内部相同，仍须一个比较论证。直接使用跳跃的指示数据会在弧端点留下未控制的极限；采用连续的短弧上界，可以在每个边界点应用已证最大值原理。

\Needspace{6\baselineskip}
\begin{lemma}[有限边界例外集的比较]\label{b3:lem:16-finite-boundary-exceptions}
设 \(v:B(0,1)\to\mathbb R\) 调和且有上界，\(E\subset\partial B(0,1)\) 为有限集。若
\[
 \limsup_{B(0,1)\ni z\to\xi}v(z)\le0
 \quad(\xi\in\partial B(0,1)\setminus E),
\]
则 \(v\le0\) 于圆盘。特别地，若 \(v\) 有界且在除 \(E\) 外的每个边界点均趋于零，则 \(v=0\)。
\end{lemma}
\begin{proof}
取 \(M\ge0\) 使 \(v\le M\)。给定 \(\eta>0\)，用有限条开弧覆盖 \(E\)，使总角长小于 \(\eta\)。在每条弧中取一个连续函数，在相应例外点的小邻弧上等于一，在大弧外等于零，并在中间线性过渡。取这些函数的最大值，得到连续圆周函数 \(\chi_\eta\)，满足
\[
 0\le\chi_\eta\le1,\quad
 \chi_\eta=1\text{ 于 }E,\quad
 \int_0^{2\pi}\chi_\eta(e^{\mathrm i t})\,\mathrm dt<\eta.
\]
当 \(E\) 为空时直接取 \(\chi_\eta=0\)。令
\(h_\eta=\mathcal P_1\chi_\eta\)，则 \(h_\eta\ge0\)，且连续取边值 \(\chi_\eta\)。

在 \(E\) 以外，\(\limsup(v-Mh_\eta)\le-M\chi_\eta\le0\)；在 \(E\) 上，由 \(v\le M\) 与 \(h_\eta\to1\)，同样得到上极限不超过零。因此调和函数 \(v-Mh_\eta\) 在有界圆盘的每个边界点均满足上极限条件，\cref{b3:cor:15-boundary-limsup}给出
\(v\le Mh_\eta\)。固定 \(z\) 后，核满足
\[
 0\le h_\eta(z)
 \le\frac{1+|z|}{1-|z|}\,\frac{\eta}{2\pi}.
\]
令 \(\eta\downarrow0\) 得 \(v(z)\le0\)。最后若 \(v\) 有界且边界极限为零，对 \(v\) 与 \(-v\) 分别应用前一结论，得到 \(v=0\)。
\end{proof}

上半平面给出这一引理的一个用途：将它映到圆盘后，无穷远只对应一个边界点。实轴上的积分核可以先直接构造。对 \(z=x+\mathrm i y\)、\(y>0\)，函数
\[
 \frac1\pi\operatorname{Im}\frac1{t-z}
 =\frac1\pi\frac{y}{(x-t)^2+y^2}
\]
关于 \(z\) 调和且为正。代入 \(t=x+ys\)，它对 \(t\) 的积分成为
\(\pi^{-1}\int_{\mathbb R}(1+s^2)^{-1}\,\mathrm ds=1\)，所以也能用来延拓有界数据。

\Needspace{6\baselineskip}
\begin{theorem}[上半平面的 Poisson 积分]\label{b3:thm:16-half-plane-poisson}
设 \(\psi:\mathbb R\to\mathbb R\) 为有界 Borel 可测函数，定义
\[
 U(x,y)=\frac1\pi\int_{\mathbb R}
       \frac{y\,\psi(t)}{(x-t)^2+y^2}\,\mathrm dt,
 \quad x\in\mathbb R,\ y>0.
\]
则 \(U\) 在上半平面调和，且 \(|U|\le\|\psi\|_\infty\)。若 \(\psi\) 在 \(x_0\in\mathbb R\) 连续，则
\[
 U(x,y)\longrightarrow\psi(x_0)
 \quad\text{当 }(x,y)\to(x_0,0),\ y>0.
\]
若 \(\psi\) 在整条实轴连续，则 \(U\) 是唯一有界、在上半平面调和并连续取实轴边值 \(\psi\) 的函数。
\end{theorem}
\begin{proof}
核为正且积分为一，因此积分绝对收敛，并给出 \(|U|\le\|\psi\|_\infty\)。固定上半平面中的紧集 \(K\)。其上的 \(x,y\) 有界，且 \(y\) 有正下界。核及其至多二阶的 \(x,y\) 偏导数，在 \(t\) 的有界区间上统一有界，在 \(|t|\) 足够大时由 \(C_K/(1+t^2)\) 控制；这一点可由核的有理表达直接求导得到。乘以有界的 \(\psi\) 后仍有共同可积控制，故可在积分号下求二阶导数。核本身是关于 \(z\) 的全纯函数的虚部，Laplace 算子为零，因而 \(\Delta U=0\)。

令 \(M=\|\psi\|_\infty\)。固定连续点 \(x_0\) 与 \(\gamma>0\)，取 \(\delta>0\)，使 \(|t-x_0|<\delta\) 时
\(|\psi(t)-\psi(x_0)|<\gamma\)。若 \(|x-x_0|<\delta/2\)，则远部满足
\(|t-x|\ge|t-x_0|/2\)，所以
\[
 \frac1\pi\int_{|t-x_0|\ge\delta}
       \frac{y}{(x-t)^2+y^2}\,\mathrm dt
 \le\frac{4y}{\pi}\int_{|s|\ge\delta}\frac{\mathrm ds}{s^2}
 =\frac{8y}{\pi\delta}.
\]
将 \(U(x,y)-\psi(x_0)\) 按这两部分拆开，近部用核的总质量一，远部用差值至多 \(2M\)，得到
\[
 |U(x,y)-\psi(x_0)|
 \le\gamma+\frac{16My}{\pi\delta}.
\]
先令 \((x,y)\to(x_0,0)\)，再令 \(\gamma\downarrow0\)，便得到任意方向的边界极限。

现在假设 \(\psi\) 在整条实轴连续，并取另一有界调和解 \(V\)。差 \(W=V-U\) 有界且连续取零实轴边值。映射
\[
 F(\zeta)=\mathrm i\,\frac{1+\zeta}{1-\zeta},
 \quad |\zeta|<1,
\]
将单位圆盘双全纯地映到上半平面，因为
\[
 \operatorname{Im}F(\zeta)=\frac{1-|\zeta|^2}{|1-\zeta|^2}>0,
 \quad
 F^{-1}(z)=\frac{z-\mathrm i}{z+\mathrm i}.
\]
除圆周点 \(1\) 外，\(F\) 连续取有限实轴值。链式法则与 Cauchy--Riemann 方程给出
\(\Delta(W\circ F)=|F'|^2\cdot(\Delta W)\circ F=0\)，故 \(W\circ F\) 是有界调和函数，并在除 \(1\) 外的所有圆周点趋于零。\cref{b3:lem:16-finite-boundary-exceptions}遂给 \(W\circ F=0\)，从而 \(V=U\)。证明没有要求在无穷远预先存在一个边值。

\end{proof}

有界性在唯一性证明中控制了映到圆周点 \(1\) 的无穷远部分。若放弃这项条件，零函数与 \(U(x,y)=y\) 都连续取零实轴边值，却在内部不同。同一证明还允许在实轴上忽略有限个边界点：若两个有界调和函数在其余每个实点的差都趋于零，拉回后的例外集只是这些点的像再加上 \(1\)，仍为有限集，故两函数相等。

\ProseParagraphBreak
实轴非紧，各连续点的边界收敛本身不保证全实轴上的一致收敛。

\ProseParagraphBreak
一些积分值可以通过唯一性计算，而不必直接积分。数据
\(\psi(t)=1/(1+t^2)\) 的有界调和延拓是
\[
 U(x,y)=\frac{1+y}{x^2+(1+y)^2}
       =-\operatorname{Im}\frac1{x+\mathrm i y+\mathrm i}.
\]
右端的极点位于下半平面，在 \(y\ge0\) 上介于零与一之间，并连续取指定边值，故由唯一性就是 Poisson 积分。对数据 \(\cos(\kappa t)\)，其中 \(\kappa\in\mathbb R\)，函数
\[
 U(x,y)=e^{-|\kappa|y}\cos(\kappa x)
\]
同样直接满足调和方程、有界性和边界条件。频率越高，向上半平面内部的衰减越快；这与圆盘中 \((r/R)^{|n|}\) 的衰减是两个可明确计算的模型。

% !TeX root = ../../Book3.tex
\section{Perron 方法}
\label{b3:sec:1602}

一般区域没有可直接写出的核。Perron 方法从所有低于边界数据的次调和函数出发；圆盘上的调和替换迫使它们的上包络满足平均等式。

\subsection{次函数}
\label{b3:subsec:160201}

称连续次调和函数 \(v:D\to\mathbb R\) 为边值 \(\varphi\) 的\term[perron-cixanshu]{Perron 次函数}[Perron subfunction]，若对每个 \(\xi\in\partial D\) 有 \(\limsup_{z\to\xi}v(z)\le\varphi(\xi)\)。本章只采用连续的次函数；这个族已足以完成连续边值的构造。

常数 \(m=\min\varphi\) 是次函数。每个次函数由次调和最大值原理满足 \(v\le M=\max\varphi\)，所以取上确界不会在内部发散。连续性允许在任何内部闭圆盘上作 Poisson 替换。

\subsection{Perron 族}
\label{b3:subsec:160202}

把全部次函数组成的族记为 \(\mathcal P_\varphi\)，上包络记为 \(H_\varphi(z)=\sup_{v\in\mathcal P_\varphi}v(z)\)。有限最大值仍在族中。

若 \(B\Subset D\)，把 \(v\) 在 \(B\) 内换成同边值的 Poisson 解，外部保持 \(v\)，记为 \(v^B\)。比较原理给 \(v^B\ge v\)。该函数在圆周连续，且仍次调和；下文通过先比较原次函数、再比较替换部分给出跨圆周的完整核查。替换只影响内部紧集，所以原边界限制不变。这个封闭性是上包络证明的核心。

\subsection{Perron 解}
\label{b3:subsec:160203}

\begin{theorem}[Perron]\label{b3:thm:16-perron}
有界区域 \(D\) 与连续边值 \(\varphi\) 的包络 \(H_\varphi\) 在 \(D\) 调和，并满足 \(m\le H_\varphi\le M\)。
\end{theorem}
\begin{proof}
固定内部圆盘 \(B\) 及 \(a\in B\)。选 \(v_n\) 使 \(v_n(a)\to H_\varphi(a)\)，逐次取有限最大值，得到递增序列，再在 \(B\) 调和替换。所得 \(h_n\) 递增、有界，在 \(B\) 局部一致收敛到调和函数 \(h\)，且 \(h(a)=H_\varphi(a)\)、\(h\le H_\varphi\)。
任取另一项 \(v\in\mathcal P_\varphi\)，将 \(\max(v,h_n)\)（在全域使用对应替换后的函数）再提升，得到调和函数 \(k_n\ge h_n,v\)。它们同样递增有界，极限 \(k\) 调和。中心满足 \(h(a)\le k(a)\le H_\varphi(a)=h(a)\)，所以非负调和函数 \(k-h\) 在内部为零，由强最小值原理恒零。于是 \(v\le h\) 于 \(B\)，取全族上确界得 \(H_\varphi\le h\)。因此包络在每个内部圆盘调和。
\end{proof}
证明没有先断言任意上包络连续；连续性由局部调和识别产生。这一步避免了把上确界与上半连续正则化混同。

\subsection{唯一性}
\label{b3:subsec:160204}

若一个经典解 \(u\) 已存在，则它本身是次函数，每个次函数又由比较原理不超过 \(u\)，故 \(H_\varphi=u\)。因此 Perron 构造不产生第二种不同的经典解。

没有边界正则性时，\(H_\varphi\) 仍是内部调和函数，却可能忽略某些指定的边界值。去心圆盘上，外圆取零、孤立边界点原点取一，得到 Perron 解恒零：有界调和函数的孤立奇点可去，延拓后外圆边值迫使零函数。该例将在下一节同时说明障碍函数为何必要。


\ProseParagraphBreak

调和替换在圆周上的粘合可以作如下严格核查。设 \(w=v^B\)，并在测试圆盘 \(Q\Subset D\) 取调和函数 \(h\)，满足 \(w\le h\) 于 \(\partial Q\)。由于 \(v\le w\)，先由 \(v\) 的次调和比较得到 \(v\le h\) 于 \(Q\)。在 \(Q\cap B\) 中，\(w-h\) 调和；其边界分成 \(\partial Q\cap\overline B\) 与 \(Q\cap\partial B\)，前者由假设非正，后者因 \(w=v\le h\) 非正。最大值原理便给 \(w\le h\) 于交集。其余部分 \(w=v\le h\)，所以整个 \(Q\) 比较成立。这里先比较原次函数，再比较替换部分，避免在交界处循环使用尚未证明的次调和性。

\ProseParagraphBreak

Perron 证明中选择序列只需要在一个点逼近包络，而不要求全域同时选出最优函数。具体可令 \(v_n(a)>H_\varphi(a)-1/n\)，将 \(m,v_1,\ldots,v_n\) 的最大值记为 \(s_n\)，再替换得 \(w_n=s_n^B\)。序列 \(w_n\) 在全域递增，且属于同一 Perron 族；其圆盘限制 \(h_n\) 调和。由于 \(h_n(a)\ge s_n(a)\) 且 \(h_n\le H_\varphi\)，中心极限正是包络。对任意竞争者 \(v\)，使用 \(\max(v,w_n)\) 而非只定义在 \(B\) 中的 \(\max(v,h_n)\)，就得到合法的全域次函数。这些域的区别保证第二轮替换仍保留边界约束。

\ProseParagraphBreak

取上确界并不通常与取极限交换。Perron 证明采用的事实是：对每个固定竞争者 \(v\)，新调和极限 \(k\) 在圆盘上支配它，同时 \(k(a)=h(a)\)；非负调和函数的强最小值原理强迫 \(k=h\)。于是一个最初只在中心被优化的序列，最终在整个圆盘支配全部竞争者。这个步骤使用线性的调和方程和正性，才产生包络的调和性；只看有限最大值的次调和封闭性还不足以完成证明。

\ProseParagraphBreak

去心圆盘例子的包络恒零也能在不预先调用可去奇点定理的情况下证明。对任意次函数 \(v\)，外圆边界上限零，内边界点上限至多一，故 \(v\le1\)。给定 \(\varepsilon>0\)，在足够小的内圆 \(|z|=r\) 上 \(v\le1+\varepsilon\)，环域上的调和函数
\[
q_r(z)=(1+\varepsilon)\frac{\log|z|}{\log r}
\]
外圆值零、内圆值 \(1+\varepsilon\)。比较给 \(v(z)\le q_r(z)\)；固定 \(z\ne0\) 让 \(r\downarrow0\)，得到 \(v(z)\le0\)。常数零又是合法次函数，所以包络恰为零。孤立边值一在内部没有影响，这一计算完全在当前已证明的环域比较范围内进行。

\ProseParagraphBreak

边界数据的单调性不需要域正则：若 \(\varphi\le\psi\)，则 \(\mathcal P_\varphi\subseteq\mathcal P_\psi\)，故 \(H_\varphi\le H_\psi\)。加常数与乘正数也分别对应次函数族的平移与缩放，所以 \(H_{\varphi+c}=H_\varphi+c\)、\(H_{a\varphi}=aH_\varphi\) 对 \(a>0\) 成立。两者推出
\[
\|H_\varphi-H_\psi\|_{L^\infty(D)}\le\|\varphi-\psi\|_{L^\infty(\partial D)}.
\]
一般加法的线性不能只由族包含轻率推出；本章在正则域用经典解的唯一性证明完整线性，已经足够用于测度表示。

\ProseParagraphBreak

圆环提供检验一般存在理论的另一组显式解。设 \(r_0<|z|<R\)，令 \(x=\log(|z|/r_0)\)、\(L=\log(R/r_0)\)，则 \(0<x<L\)。极坐标下 Laplace 方程乘以 \(|z|^2\) 后成为 \(u_{xx}+u_{\theta\theta}=0\)。对内、外边界各一个 Fourier 模式，解分别为
\[
 \frac{\sinh(n(L-x))}{\sinh(nL)}e^{\mathrm{i}n\theta},\qquad
 \frac{\sinh(nx)}{\sinh(nL)}e^{\mathrm{i}n\theta}\quad(n\ge1).
\]
它们在相应边界取一倍模式，在另一边界取零；对负频率取共轭。零频率则为 \(1-x/L\) 与 \(x/L\)。每一式都可直接求导检验，唯一性保证它与 Perron 解相同。
\ProseParagraphBreak
若内、外边界为三角多项式，有限叠加便得到解。对一般连续周期数据，可由非负的 Fej\'er 核逼近：
\[
 F_N(t)=\frac1N\left|\sum_{j=0}^{N-1}e^{\mathrm{i}jt}\right|^2,
 \qquad \frac1{2\pi}\int_{-\pi}^{\pi}F_N(t)\,\mathrm{d}t=1.
\]
在 \(|t|\ge\delta\) 上，有限几何和给 \(F_N(t)\le C_\delta/N\)；在小弧上用数据的一致连续性。因此卷积 \(F_N*\varphi\) 一致趋于 \(\varphi\)，每次卷积都是三角多项式。对这些数据的圆环解，最大值原理使内部一致差不超过两条边界数据的最大差，故在闭圆环上一致收敛。极限在内部调和，因为任意小圆盘上的平均值性质可通过一致极限；边值也正确。这给出圆环上存在性的第二个完整构造，并显示对数项恰来自角向零频率。

% !TeX root = ../../Book3.tex
\section{边界正则性}
\label{b3:sec:1603}

Perron 包络已经在内部调和，但这还没有说明它是否取到每一个指定的边界值。本节把问题集中到一个边界点：若能找到在这个点趋零、远离这个点却保持正下界的超调和函数，就能用它同时控制所有连续边值的误差。这样的函数称为障碍；它将局部几何信息转化为边界比较。

\subsection{障碍函数}
\label{b3:subsec:160301}

在边界点 \(\xi\) 附近，我们希望一个正函数只在趋向 \(\xi\) 时变小。仅要求它在域内逐点为正并不够：沿另一个边界点趋近时，它仍可能趋于零。定义中的环带下界排除这种情形，也为稍后的截断和拼接留出余量。

\Needspace{6\baselineskip}
\begin{definition}\label{b3:def:16-barrier}
设 \(D\subseteq\mathbb C\) 为有界区域，\(\xi\in\partial D\)。若存在 \(r_0>0\) 及连续函数
\[
 b:D\cap B(\xi,r_0)\longrightarrow(0,\infty),
\]
使 \(b\) 在该开集的每个连通分支超调和，且满足
\[
 \lim_{z\to\xi,\ z\in D}b(z)=0,
 \quad
 \inf_{\substack{z\in D\\\delta\le|z-\xi|\le r'}}b(z)>0
 \quad(0<\delta<r'<r_0),
\]
则称 \(b\) 为 \(\xi\) 处的连续正超调和\term[zhangai-hanshu]{障碍函数}[barrier]，简称局部障碍。环带与 \(D\) 的交集为空时，下界条件不另加限制。
\end{definition}
这里的超调和是指 \(-b\) 次调和。函数 \(b\) 不必连续延拓到整个局部边界；在一个可能不闭的集合 \(D\cap\{|z-\xi|=s\}\) 上逐点为正，也不能单凭连续性推出统一正下界。因而定义直接要求内部环带上的统一下界，而不以人工圆周的紧性代替它。

\Needspace{6\baselineskip}
\begin{lemma}[局部障碍的全域化]\label{b3:lem:16-global-barrier}
设 \(D\subseteq\mathbb C\) 为有界区域，\(\xi\in\partial D\)，且 \(\xi\) 有连续正超调和局部障碍 \(b\)。则存在有界、连续、正的超调和函数 \(B\) 于 \(D\)，在 \(\xi\) 的某个域内邻域等于 \(b\)，并满足
\[
 \lim_{z\to\xi,\ z\in D}B(z)=0,
 \quad
 \inf_{D\setminus B(\xi,\delta)}B>0
 \quad(\delta>0).
\]
最后一式只要求于非空的集合。
\end{lemma}
\begin{proof}
设 \(b\) 定义在 \(D\cap B(\xi,r_0)\)。取 \(0<s<r_0/2\)。由环带条件，可选 \(c>0\)，使
\[
 b(z)\ge2c
 \quad\left(z\in D,\ s\le|z-\xi|\le2s\right).
\]
若该环带中没有域内点，任取 \(c>0\) 即可。定义
\[
 B(z)=
 \begin{cases}
 \min\{b(z),c\},&|z-\xi|<3s/2,\\
 c,&|z-\xi|\ge3s/2.
 \end{cases}
\]
在整个域内环带 \(s<|z-\xi|<2s\) 上，两种表达都等于 \(c\)。因此拼接处已有一个共同的常值邻域，而不是仅在人工圆周上取值相同。函数 \(B\) 连续且处处为正。有限个超调和函数的最小值仍超调和，这是\cref{b3:prop:15-subharmonic-operations}对其负号取最大值的结论；再用\cref{b3:prop:15-local-comparison}的局部性质，得到全域超调和性。

当 \(z\to\xi\) 时 \(b(z)\to0\)，所以在充分小的邻域有 \(b<c\)，从而 \(B=b\)。对远处的下界，若 \(\delta\ge s\)，则 \(D\setminus B(\xi,\delta)\) 上 \(B=c\)。若 \(0<\delta<s\)，原障碍在 \(\delta\le|z-\xi|\le s\) 上有正下界；将它与 \(c\) 取较小值，便控制了整个 \(D\setminus B(\xi,\delta)\)。
\end{proof}
圆盘给出一个可以直接核验的全域障碍。设 \(D=B(a,R)\)、\(|\xi-a|=R\)，令
\[
 b(z)=1-\operatorname{Re}\frac{(z-a)\overline{(\xi-a)}}{R^2}.
\]
它是调和函数，而且
\[
 2R^2b(z)=R^2-|z-a|^2+|z-\xi|^2\ge|z-\xi|^2.
\]
在圆盘内部第一项严格为正，因此 \(b>0\)；趋于 \(\xi\) 时 \(b\to0\)，而 \(|z-\xi|\ge\delta\) 时 \(b(z)\ge\delta^2/(2R^2)\)。这个恒等式把“远离指定边界点”需要的统一下界直接算了出来。

\subsection{正则边界点}
\label{b3:subsec:160302}

边界正则性首先是解对数据的性质，不依赖我们是否已经找到某种几何障碍。由\cref{b3:thm:16-perron}，有界区域上每份连续实边值都有一个内部调和的 Perron 包络；定义只问这个包络在一个指定边界点的极限。

\Needspace{6\baselineskip}
\begin{definition}\label{b3:def:16-regular-point}
设 \(D\subseteq\mathbb C\) 为有界区域，\(\xi\in\partial D\)。对每个 \(\varphi\in C(\partial D;\mathbb R)\)，令 \(H_\varphi\) 为由连续次调和 Perron 次函数组成的上包络。若
\[
 \lim_{z\to\xi,\ z\in D}H_\varphi(z)=\varphi(\xi)
 \quad\left(\varphi\in C(\partial D;\mathbb R)\right),
\]
则称 \(\xi\) 为 \(D\) 的\term[zhengze-bianjiedian]{正则边界点}[regular boundary point]。若 \(\partial D\) 的每个点都正则，称 \(D\) 为\term[zhengze-yu]{正则域}[regular domain]。
\end{definition}
量词“每个连续边值”不能省略。常数数据在任何边界点都会取到正确极限，不能用一份常数数据检验正则性。下面的等价命题既给充分条件，也说明局部障碍确实描述了这一性质。

\Needspace{6\baselineskip}
\begin{proposition}[障碍与边界正则性]\label{b3:prop:16-barrier-regular}
设 \(D\subseteq\mathbb C\) 为有界区域，\(\xi\in\partial D\)。则 \(\xi\) 为正则边界点，当且仅当它具有连续正超调和局部障碍。若 \(\xi\) 正则，还可以取一个在整个 \(D\) 调和的正障碍，且它在每个 \(D\setminus B(\xi,\delta)\) 上有统一正下界。
\end{proposition}
\begin{proof}
先设局部障碍存在，由\cref{b3:lem:16-global-barrier}取全域障碍 \(B\)。固定 \(\varphi\in C(\partial D;\mathbb R)\) 和 \(\varepsilon>0\)。选 \(\delta>0\)，使
\[
 |\varphi(\eta)-\varphi(\xi)|<\varepsilon
 \quad\left(\eta\in\partial D,\ |\eta-\xi|<\delta\right).
\]
设 \(M=\max_{\partial D}|\varphi-\varphi(\xi)|\)。全域障碍在 \(D\setminus B(\xi,\delta/2)\) 上有正下界 \(c_\delta\)，取 \(A\ge0\) 使 \(Ac_\delta\ge M\)。若该集合为空，所有边界点都在上述邻域中，取 \(A=0\) 即可。

函数
\[
 L(z)=\varphi(\xi)-\varepsilon-AB(z),
 \quad U(z)=\varphi(\xi)+\varepsilon+AB(z)
\]
分别连续次调和、连续超调和。在距离 \(\xi\) 小于 \(\delta\) 的边界点，正性和数据连续性给 \(\limsup L\le\varphi\)、\(\liminf U\ge\varphi\)。在其余边界点，域内的趋近点最终满足 \(|z-\xi|\ge\delta/2\)，统一下界 \(c_\delta\) 及 \(Ac_\delta\ge M\) 给出相同结论。因此 \(L\) 属于 Perron 次函数族。对任意该族中的 \(v\)，函数 \(v-U\) 次调和，边界上极限非正；\cref{b3:prop:15-subharmonic-maximum}在有界域上给 \(v\le U\)。取上确界得到
\[
 \varphi(\xi)-\varepsilon-AB(z)
 \le H_\varphi(z)\le
 \varphi(\xi)+\varepsilon+AB(z).
\]
先令 \(z\to\xi\)，利用 \(B(z)\to0\)，再令 \(\varepsilon\downarrow0\)，便得正则性。

反过来，只假定 \(\xi\) 正则。取连续边值 \(\psi(\eta)=|\eta-\xi|^2\)，并令 \(B=H_\psi\)。函数 \(v(z)=|z-\xi|^2\) 连续次调和，且在每个边界点取到 \(\psi\)，所以它本身就是合法次函数。因此
\[
 0<|z-\xi|^2\le B(z)\le\max_{\partial D}|\eta-\xi|^2
 \quad(z\in D).
\]
Perron 定理保证 \(B\) 调和；\(\xi\) 的正则性又给 \(B(z)\to\psi(\xi)=0\)。而 \(|z-\xi|\ge\delta\) 时 \(B(z)\ge\delta^2\)，故 \(B\) 是所需全域正调和障碍。这个构造没有使用其他边界点的正则性。
\end{proof}
因此只需检验一份与指定点相关的数据：若边值 \(|\eta-\xi|^2\) 的 Perron 包络在 \(\xi\) 趋于零，上面的构造就给出障碍，继而迫使所有连续数据在 \(\xi\) 取到正确极限。正则性也因而是局部的：局部障碍在相同的域内边界几何上可以直接使用，全域部分由截断补齐。

\ProseParagraphBreak

在正则域上，\(H_\varphi\) 在每个边界点取连续数据的值，连同内部的调和性便给 \(H_\varphi\in C(\overline D)\)。经典解的唯一性来自最大值原理。因此，Perron 法的内部存在性与障碍的边界到达性质合在一起，完成了正则域上连续 Dirichlet 问题的存在唯一性。

\subsection{外锥条件}
\label{b3:subsec:160303}

外锥属于补集。具体地，若 \(\mathbb C\setminus D\) 在 \(\xi\) 附近包含一个以 \(\xi\) 为顶点、角宽为正的实心扇形，域的局部部分便落在余下角宽 \(\beta<2\pi\) 的扇形中。\(\beta\) 表示容纳域的角宽；它不是补集外锥的角宽。选定旋转角 \(\theta_0\) 后，把域内点写成
\[
 z=\xi+r e^{\mathrm{i}(\theta_0+\theta)},
 \quad 0<r<r_0,\quad |\theta|<\beta/2.
\]
因为 \(\beta<2\pi\)，这里可以选取一个连续辐角分支。域只须包含在这个扇形中，不必等于扇形，也不要求原边界是直线。

\ProseParagraphBreak

极坐标下的调和方程提示取 \(r^\alpha\cos(\alpha\theta)\)：径向二阶项和角向二阶项恰好相消。但障碍还要求整个允许角范围内保持正下界，所以 \(\alpha\) 的选择不仅由调和方程决定。

\Needspace{6\baselineskip}
\begin{proposition}[外锥条件与正则性]\label{b3:prop:16-exterior-cone-regular}
设 \(D\subseteq\mathbb C\) 为有界区域，\(\xi\in\partial D\)。若存在 \(r_0>0\)、\(\theta_0\in\mathbb R\) 及 \(0<\beta<2\pi\)，使
\[
 D\cap B(\xi,r_0)
 \subseteq\{\,\xi+r e^{\mathrm{i}(\theta_0+\theta)}
       \mid 0<r<r_0,\ |\theta|<\beta/2\,\},
\]
则 \(\xi\) 正则。对任意 \(0<\alpha<\pi/\beta\)，函数
\[
 b(z)=r^\alpha\cos(\alpha\theta)
\]
都是 \(\xi\) 处的连续正调和局部障碍。
\end{proposition}
\begin{proof}
在上述辐角分支上，\(b\) 是 \(\bigl(e^{-\mathrm{i}\theta_0}(z-\xi)\bigr)^\alpha\) 的实部，故调和。令 \(c_\alpha=\cos(\alpha\beta/2)\)。严格不等式 \(\alpha\beta/2<\pi/2\) 给 \(c_\alpha>0\)，并且
\[
 c_\alpha r^\alpha\le b(z)\le r^\alpha.
\]
上界给 \(z\to\xi\) 时 \(b(z)\to0\)，下界给每个 \(\delta\le r\le r'<r_0\) 上的统一正下界 \(c_\alpha\delta^\alpha\)。因此 \(b\) 满足局部障碍的全部条件，由\cref{b3:prop:16-barrier-regular}得到正则性。
\end{proof}
\begin{figure}[htbp]
\centering
\begin{tikzpicture}[>=stealth,font=\small,scale=1.12]
 \fill[gray!30] (0,0)--(120:2) arc (120:240:2)--cycle;
 \draw[dashed] (0,0) circle (2);
 \draw (120:2)--(0,0)--(240:2);
 \draw[dotted,->] (0,0)--(2.35,0) node[right] {\(\theta=0\)};
 \draw[->] (-120:0.65) arc (-120:120:0.65);
 \node[right] at (0.68,0.22) {\(\beta\)};
 \draw[thick] (-0.38,1.6)..controls (-0.4,0.7) and (-0.1,0.35)..(0,0)
   ..controls (0.3,-0.65) and (0.25,-1.3)..(0.8,-1.8);
 \node at (1.18,0.8) {\(D\)};
 \node[align=center] at (-1.3,0) {补集\\外锥};
 \fill (0,0) circle (0.025) node[below left] {\(\xi\)};
 \node[above right] at (45:2) {人工圆周 \(r=r_0\)};
\end{tikzpicture}
\caption[外锥与正障碍的扇形]{阴影锥包含于补集，域在局部只包含于余下角宽 \(\beta<2\pi\) 的扇形，不要求等于该扇形。粗曲线仅示意可能的域边界。选择 \(\alpha\beta/2<\pi/2\) 后，整个允许角范围内都有 \(\cos(\alpha\theta)\ge\cos(\alpha\beta/2)>0\)，故人工圆周上障碍具有统一正下界。}
\label{b3:fig:16-exterior-cone-barrier}
\end{figure}

\cref{b3:fig:16-exterior-cone-barrier}中的粗曲线示意原域边界，虚线圆周则是局部化人为加入的边界。余弦的统一正下界在整个环带有效，因而可以使用\cref{b3:lem:16-global-barrier}；不必假定人工圆周与域的交集是紧的。若取临界值 \(\alpha=\pi/\beta\)，余弦在两条侧边趋于零，环带上的正下界就可能消失。这个临界函数仍然调和，却未必符合当前障碍的定义。

\ProseParagraphBreak

例如局部 Lipschitz 边界可旋转成图像 \(y=g(x)\)，其中 \(g(0)=0\)、\(|g(x)-g(x')|\le L|x-x'|\)，域局部位于 \(y>g(x)\)。由于 \(g(x)\ge-L|x|\)，下方的锥 \(y<-(L+1)|x|\) 包含于补集。于是 Lipschitz 边界点满足外锥条件；光滑边界及角宽小于 \(2\pi\) 的多边形角点也属于这一范围。再入角对应 \(\beta>\pi\)，并不妨碍正则性。

\Needspace{6\baselineskip}
\begin{proposition}[外锥下的局部 Hölder 估计]\label{b3:prop:16-corner-holder}
设 \(D\subseteq\mathbb C\) 为有界区域，\(\xi\in\partial D\)。假设某个 \(r_0>0\) 下，\(D\cap B(\xi,r_0)\) 包含于以 \(\xi\) 为顶点、角宽 \(0<\beta<2\pi\) 的扇形。设 \(\varphi\in C(\partial D;\mathbb R)\)，\(H_\varphi\) 为连续次函数 Perron 包络。固定 \(0<\alpha<\pi/\beta\)。若存在 \(H\ge0\) 和 \(0<\gamma\le1\)，使
\[
 |\varphi(\eta)-\varphi(\xi)|\le H|\eta-\xi|^\gamma
 \quad\left(\eta\in\partial D\cap B(\xi,r_0)\right),
\]
则存在 \(r_1>0\)、\(C\ge0\)，使
\[
 |H_\varphi(z)-\varphi(\xi)|
 \le C|z-\xi|^{\min(\alpha,\gamma)}
 \quad\left(z\in D\cap B(\xi,r_1)\right).
\]
常数可依赖局部扇形、所选指数、半径及边值数据。
\end{proposition}
\begin{proof}
取 \(b=r^\alpha\cos(\alpha\theta)\)，令
\[
 q=\min(1,\gamma/\alpha),\quad p=\alpha q=\min(\alpha,\gamma),
 \quad c_\alpha=\cos(\alpha\beta/2)>0.
\]
在域内部 \(b>0\)、\(\Delta b=0\)，链式法则给
\[
 \Delta(b^q)=q(q-1)b^{q-2}|\nabla b|^2\le0,
 \quad c_\alpha^q r^p\le b^q\le r^p.
\]
故 \(b^q\) 是连续正超调和局部障碍。由\cref{b3:lem:16-global-barrier}取其全域化 \(B\)，并选 \(0<\rho<r_0\)，使 \(B=b^q\) 于 \(D\cap B(\xi,\rho)\)。

若边界点 \(\eta\) 满足 \(0<|\eta-\xi|<\rho/2\)，则
\[
 \liminf_{z\to\eta,\ z\in D}B(z)
 \ge c_\alpha^q|\eta-\xi|^p,
 \quad
 |\varphi(\eta)-\varphi(\xi)|
 \le H(\rho/2)^{\gamma-p}|\eta-\xi|^p.
\]
因此可用同一个倍数 \(A B\) 支配这部分边界误差。在其余边界点，域内趋近点最终离 \(\xi\) 至少 \(\rho/4\)，而 \(B\) 在这个范围有统一正下界；增大 \(A\) 即可支配其余有界误差。在 \(\xi\) 本身，\(B\to0\)，误差也为零。

于是 \(\varphi(\xi)-AB\) 是 Perron 次函数，\(\varphi(\xi)+AB\) 是支配所有次函数的超调和上界。和障碍正则性证明中的比较相同，得到
\[
 |H_\varphi(z)-\varphi(\xi)|\le AB(z)\le A|z-\xi|^p
 \quad\left(z\in D\cap B(\xi,\rho)\right).
\]
取 \(C=A\)、\(r_1=\rho\) 即得结论。
\end{proof}
当边值指数 \(\gamma\) 小于障碍指数 \(\alpha\) 时，取 \(q=\gamma/\alpha<1\) 会使障碍在顶点附近下降得慢一些，足以控制较粗糙的边值误差；凹幂产生的超调和性保留了比较方向。仅有某个障碍存在，并不保证它具有这样的幂次上下界，所以定性的边界到达不能自动加强为固定指数的 Hölder 估计。

\ProseParagraphBreak

再入角确实可能降低解的边界光滑性。在截断扇形 \(0<r<1\)、\(|\theta|<\beta/2\) 上，取 \(\pi<\beta<2\pi\)，则
\[
 u(r,\theta)=r^{\pi/\beta}\cos(\pi\theta/\beta)
\]
调和、连续到闭域，并在两条侧边及顶点取零。在中轴上，\(u(r,0)/r=r^{\pi/\beta-1}\to\infty\)，所以它在顶点不是 Lipschitz 函数，尽管靠近顶点的边界数据恒为零。正则性保证连续到达，并未保证一阶边界正则性。

\subsection{障碍算例与 Wiener 判据概览}
\label{b3:subsec:160304}

外接圆提供另一种直接可算的障碍，但须限制接触集。若同一个外圆沿一段边界接触，下面的对数会同时在这些点趋于零，不能单独识别指定点。

\Needspace{6\baselineskip}
\begin{proposition}[唯一接触的外球障碍]\label{b3:prop:16-exterior-ball-barrier}
设 \(D\subseteq\mathbb C\) 为有界区域，\(\xi\in\partial D\)。若存在 \(a\in\mathbb C\)、\(R>0\)、\(r_0>0\)，使
\[
 \overline{B(a,R)}\subseteq\mathbb C\setminus D,
 \quad |\xi-a|=R,
 \quad
 \overline D\cap\overline{B(\xi,r_0)}\cap\overline{B(a,R)}=\{\xi\},
\]
则 \(b(z)=\log(|z-a|/R)\) 在 \(D\cap B(\xi,r_0)\) 是连续正调和障碍，因而 \(\xi\) 正则。
\end{proposition}
\begin{proof}
域内有 \(|z-a|>R\)，故 \(b>0\)，且它在不含 \(a\) 的该开集调和。趋于 \(\xi\) 时，由 \(|\xi-a|=R\) 得 \(b\to0\)。固定 \(0<\delta<r'<r_0\)，考虑紧集
\[
 F=\overline D\cap\{\,z\mid\delta\le|z-\xi|\le r'\,\}.
\]
若 \(F\) 非空，唯一接触条件使 \(F\) 与 \(\overline{B(a,R)}\) 不交。连续函数 \(|z-a|-R\) 因而在 \(F\) 上有严格正的最小值，所以 \(\log(|z-a|/R)\) 在相应域内环带上有统一正下界。由障碍判据得结论。
\end{proof}
对于 \(C^2\) 图像边界，可把接触点移到原点、切线转为水平线，写 \(D\) 局部在 \(y>g(x)\)，其中 \(g(0)=g'(0)=0\)、\(|g(x)|\le Cx^2\)。取下方圆心 \((0,-a)\) 和足够小的半径 \(a>0\)，使整个圆盘留在图像邻域且 \(1/(2a)>C\)。该圆的上半边满足
\[
 -a+\sqrt{a^2-x^2}\le-\frac{x^2}{2a}<-Cx^2\le g(x)
 \quad(0<|x|<a).
\]
圆盘位于域外，并只在原点接触。这说明光滑边界为何常能使用外球障碍；只有 Lipschitz 正则性时，则直接使用外锥更方便。

\ProseParagraphBreak

圆环的两条边界圆均正则。\subsecref{b3:subsec:160204}中的显式解表明，内圈取一、外圈取零时
\[
 u_r(z)=\frac{\log(R/|z|)}{\log(R/r)}
 \quad(r<|z|<R).
\]
固定非零内点后，\(r\downarrow0\) 时该值按 \(1/\log(R/r)\) 的量级趋于零。若把内圈缩成孤立点并在那里指定值一，\hyperref[b3:exmp:16-punctured-disk]{上一节的穿孔圆盘算例}已经用圆环比较证明 Perron 包络恒零。该点不正则，因而也不存在本节定义的障碍。

\ProseParagraphBreak

外锥和孤点的差别还可以用能量度量。下面只建立所需的相对容量及两个直接计算，为理解更一般的边界判据提供参照。

\Needspace{6\baselineskip}
\begin{definition}\label{b3:def:16-relative-capacity}
设 \(U\subseteq\mathbb C\) 为有界开集，\(K\Subset U\) 为紧集。定义 \(K\) 相对于 \(U\) 的\term[xiangdui-rongliang]{相对容量}[relative capacity]为
\[
 \operatorname{cap}(K,U)\coloneqq
 \inf\left\{\,\int_U|\nabla v|^2\,\mathrm{d}A
 \ \middle|\ v\in C_c^\infty(U;\mathbb R),\quad
 v\ge1\text{ 于 }K\text{ 的某邻域}\,\right\}.
\]
\end{definition}
\symidx{\(\operatorname{cap}(K,U)\)：紧集 \(K\) 相对于开集 \(U\) 的容量}

紧集增大时，试函数限制增加，容量不减；外域增大时，可用试函数增加，容量不增。二维能量还具有尺度不变性：换元 \(z=\xi+s\zeta\) 使梯度的平方乘 \(s^{-2}\)，面积元乘 \(s^2\)，两者抵消。这解释了为什么平面上的比较必须保留内外尺度之比。

\Needspace{6\baselineskip}
\begin{proposition}[圆盘的相对容量]\label{b3:prop:16-annulus-capacity}
若 \(0<r<R\)，则
\[
 \operatorname{cap}\bigl(\overline{B(0,r)},B(0,R)\bigr)
 =\frac{2\pi}{\log(R/r)},
 \quad
 \operatorname{cap}\bigl(\{0\},B(0,R)\bigr)=0.
\]
\end{proposition}
\begin{proof}
对任意合法试函数 \(v\)，有 \(v(re^{\mathrm{i}\theta})\ge1\)，而靠近 \(R\) 的圆周值为零。沿每条半径积分，再用带权 Cauchy--Schwarz 不等式，得到
\[
 \begin{aligned}
 1&\le\left|\int_r^R\partial_s v(se^{\mathrm{i}\theta})\,\mathrm{d}s\right|,\\
 1&\le\left(\int_r^R|\partial_s v(se^{\mathrm{i}\theta})|^2s\,\mathrm{d}s\right)
           \log(R/r).
 \end{aligned}
\]
对 \(\theta\) 积分，并用 \(|\nabla v|^2\ge|\partial_s v|^2\)，得到能量下界 \(2\pi/\log(R/r)\)。

为达到这个下界的极限，先取 \(r<r'<R'<R\)，令
\[
 w(z)=
 \begin{cases}
 1,&|z|\le r',\\
 \dfrac{\log(R'/|z|)}{\log(R'/r')},&r'<|z|<R',\\
 0,&|z|\ge R'.
 \end{cases}
\]
这个函数连续、分段光滑；沿两条连接圆周积分分部时函数值项相消，所以其弱梯度就是各片的通常梯度。直接计算得
\[
 \int_{\mathbb C}|\nabla w|^2\,\mathrm{d}A
 =\frac{2\pi}{\log^2(R'/r')}\int_{r'}^{R'}\frac{\mathrm{d}s}{s}
 =\frac{2\pi}{\log(R'/r')}.
\]
取非负光滑单位核 \(\eta_\varepsilon\)，支集半径为 \(\varepsilon\)，其中
\(\varepsilon<\tfrac12\min(r'-r,R-R')\)。磨光 \(w_\varepsilon=w*\eta_\varepsilon\) 在 \(|z|<r'-\varepsilon\) 等于一，支集包含于 \(\overline{B(0,R'+\varepsilon)}\Subset B(0,R)\)，所以属于容量的试函数类。又由弱梯度卷积公式和 Jensen 不等式，
\[
 |\nabla w_\varepsilon|^2
 \le |\nabla w|^2*\eta_\varepsilon,
 \quad
 \int|\nabla w_\varepsilon|^2\,\mathrm{d}A
 \le\int|\nabla w|^2\,\mathrm{d}A.
\]
令 \(r'\downarrow r\)、\(R'\uparrow R\)，便得相反的容量上界。

单点的情形也有明确试函数。取 \(0<\delta<R/8\)，在上述构造中令 \(r'=2\delta\)、\(R'=R/2\)，再用支集半径小于 \(\delta/2\) 的非负单位核磨光。所得函数在 \(|z|<\delta\) 等于一，支集包含于 \(B(0,3R/4)\)，能量不超过
\[
 \frac{2\pi}{\log(R/(4\delta))}\longrightarrow0
 \quad(\delta\downarrow0).
\]
容量非负，故单点的相对容量为零。
\end{proof}
圆环解 \(u_r\) 正是上述截断对数在过渡环带内的调和部分；其能量的分母与小孔影响的分母相同。容量为零表示可以用能量任意小的光滑函数覆盖这个点，并不表示它与域的欧氏距离较大。

\ProseParagraphBreak

外锥则在每个足够小的尺度留下统一的容量下界。设补集在 \(\xi\) 附近包含一个固定角宽的锥。由缩放，可取只依赖角宽的 \(0<\kappa<1/4\)，使每个充分小的 \(s>0\) 都有一个闭圆盘 \(\overline{B(a_s,\kappa s)}\) 位于该锥内，且 \(|a_s-\xi|=s/2\)。于是该圆盘包含于
\(K_s=(\mathbb C\setminus D)\cap\overline{B(\xi,s)}\)，而 \(B(\xi,2s)\subset B(a_s,3s)\)。由容量的两种单调性及圆盘公式，
\[
 \operatorname{cap}\bigl(K_s,B(\xi,2s)\bigr)
 \ge\operatorname{cap}\bigl(\overline{B(a_s,\kappa s)},B(a_s,3s)\bigr)
 =\frac{2\pi}{\log(3/\kappa)}>0.
\]
下界与 \(s\) 无关。外锥在不断缩小的尺度上保留了能量影响，孤点则没有这种性质。

\ProseParagraphBreak

\term[wiener-panju]{Wiener 判据}[Wiener criterion]把一般边界点的正则性与补集在逐层尺度上的容量条件联系起来。圆盘和外锥提供了这种尺度比较的可算模型。推广到任意闭补集时，需要用容量构造势、累积不同尺度的影响，并建立相应条件的必要性。Ransford 的《Potential Theory in the Complex Plane》\cite{Ransford1995Potential}可作为继续学习的资料。

% !TeX root = ../../Book3.tex
\section{调和测度}
\label{b3:sec:1604}

正则域上的解算子保持线性与正性。测度表示把整个解算子压缩成一族边界概率测度，共形映射再将这些测度从一个区域搬到另一个区域。

\subsection{解的表示}
\label{b3:subsec:160401}

设 \(D\) 是有界正则域。固定 \(z\in D\)，映射 \(\varphi\mapsto H_\varphi(z)\) 为 \(C(\partial D)\) 上的正线性泛函：两解相加仍给相加边值，唯一性使解算子线性；非负边值由最大值原理给非负解。常数一的解为一。紧空间上的正泛函表示遂给唯一概率测度 \(\omega_D^z\)，满足
\[
H_\varphi(z)=\int_{\partial D}\varphi\,\mathrm{d}\omega_D^z.
\]
这里使用第一卷定理11.4.1的 Riesz--Markov 表示。对于不同 \(z\)，这些测度通常互不相同；Harnack 估计使其对非负连续测试函数的积分在内部紧集上相互控制。

\subsection{共形不变性}
\label{b3:subsec:160402}

\begin{proposition}\label{b3:prop:16-harmonic-measure-conformal}
若 \(F:D\to D'\) 共形双射，且延拓为闭包间同胚，两域均有界正则，则
\(F_\#\omega_D^z=\omega_{D'}^{F(z)}\)。
\end{proposition}
\begin{proof}
对 \(\varphi\in C(\partial D')\)，\(H_\varphi\circ F\) 调和，因为 \(\Delta(u\circ F)=|F'|^2\cdot(\Delta u)\circ F\)。它连续取边值 \(\varphi\circ F\)，故由唯一性等于 \(H_{\varphi\circ F}\)。两边在 \(z\) 的值给两测度对全部连续测试函数积分相等，表示的唯一性得结论。
\end{proof}
若映射只在内部存在，边界像可能需要素端而不是通常拓扑边界。此时不能无说明地写相同的推前公式。

\subsection{概率解释的预告}
\label{b3:subsec:160403}

概率论中，调和测度可解释为平面 Brownian 运动首次退出区域时的位置分布。这个解释需要先构造过程，再证明退出时间、停时和边界到达的性质；它不是概率测度表示定理本身的内容。

圆盘中心的退出分布由旋转对称性应当均匀，与 Poisson 核的中心值一致。非中心处则偏向附近边界弧。本章使用的是已构造的解与测度，不以这一直观概率解释代替 Perron 证明。


\ProseParagraphBreak

解表示蕴含测度之间的定量关系。取 \(z,w\) 属于同一内部紧集，用有限 Harnack 链得到 \(C^{-1}H_\varphi(w)\le H_\varphi(z)\le C H_\varphi(w)\) 对所有非负连续 \(\varphi\) 成立。正测度的序关系因此给
\[
C^{-1}\omega_D^w(E)\le\omega_D^z(E)\le C\omega_D^w(E)
\]
对每个 Borel 边界集合 \(E\) 成立：先由连续函数积分不等式得到测度支配，再用正则性从紧集、开集推广。尤其是否具有零调和测度与内部观察点无关，虽然具体权重随点而变。这个结论只比较同一区域内的点；跨域比较需要域的关系或共形映射。

\ProseParagraphBreak

对边界弧 \(I=[\alpha,\beta]\) 的圆盘权重，中心值为 \((\beta-\alpha)/(2\pi)\)。非中心处可积分核，或用 Möbius 变换 \(T_z(\zeta)=(\zeta-z)/(1-\overline z\zeta)\) 将 \(z\) 搬到原点，得到 \(\omega^z(I)=|T_z(I)|/(2\pi)\)，这里弧长在像的单位圆上测量。原来的近弧会被展开，远弧会被压缩；这把 Poisson 核分母 \(|e^{\mathrm{i}t}-z|^2\) 中的距离效应转成圆周变换的导数。必须取相应连续像弧，不能只减两个主辐角而忽略跨越分支切线的情形。

\ProseParagraphBreak

例如上半平面的区间 \([a,b]\) 在 \(x+\mathrm{i}y\) 处的权重是
\[
\omega^{x+\mathrm{i}y}([a,b])=\frac1\pi\left(\arctan\frac{b-x}{y}-\arctan\frac{a-x}{y}\right).
\]
取 \(y\downarrow0\)，当 \(x\in(a,b)\) 时极限一，区间外极限零，端点径向极限二分之一。这个算式既检验归一化，也显示端点与弧内部的边界行为不同。若同时改变 \(x\) 和 \(y\) 以斜向接近端点，比例 \((a-x)/y\) 可趋于不同实数，因而极限也相应变化。

\ProseParagraphBreak

调和测度一般不等于普通边界长度；即使圆盘内也只有中心处是均匀的。对有尖角的域，共形变换在边界的伸缩会加强某些位置的权重，正则域这一性质本身没有承诺与弧长的密度存在。下一章光滑边界的 Green 函数法向导数将给具体密度公式，其额外边界光滑性正是微分边界核所需的条件。

% 本章小结（不计入编号）
\begin{chaptersummary}
Poisson 核以正性和质量集中同时解决圆盘存在性与边界极限。Perron 上包络通过内部调和提升成为调和函数，边界障碍再把它识别为经典解。正则域上的唯一性使解算子保持线性，正性使其具有概率测度表示。孤立缺口与外锥分别展示边界几何造成的失败和成功。
\end{chaptersummary}


\BookExercises

% 本题组自拟；各题新增训练动作见题意，未复制第三方原题。
\begin{exercise}\label{b3:ex:16-draft-01}
求单位圆周边值 \(1+2\cos t-3\sin(2t)\) 的延拓及中心值。
\end{exercise}
\begin{solution}
延拓为 \(1+2r\cos\theta-3r^2\sin(2\theta)\)，中心值一。各项为全纯多项式实部或虚部，边界值直接核对，唯一性保证结果。
\end{solution}

\begin{exercise}\label{b3:ex:16-draft-02}
若连续边值 \(\varphi\ge0\) 在一段弧上严格为正，证明其圆盘延拓在每个内部点严格为正。
\end{exercise}
\begin{solution}
连续性给一段正长度小弧及常数 \(c>0\)，使边值至少为 \(c\)。Poisson 核在该弧严格正，积分的这部分严格正；其余部分非负。
\end{solution}

\begin{exercise}\label{b3:ex:16-draft-03}
求单位圆盘中心处半圆弧的调和测度，并解释为何对一般内部点不必等于二分之一。
\end{exercise}
\begin{solution}
中心处为归一化弧长，所以等于二分之一。趋近半圆弧内部的边界点时，相应 Poisson 积分趋一，趋近另一半圆弧内部时趋零，故该权重不是常数。
\end{solution}

\begin{exercise}\label{b3:ex:16-draft-04}
证明有界调和函数在圆盘的孤立缺口可去。
\end{exercise}
\begin{solution}
在半径 \(R\) 圆周取同边值 Poisson 解 \(h\)，差 \(w\) 有界且外圆为零。在 \(\varepsilon<|z|<R\) 用调和函数 \(C\log(R/|z|)/\log(R/\varepsilon)\) 上下控制 \(w\)。固定 \(z\) 后令 \(\varepsilon\downarrow0\)，得 \(w=0\)，故由 \(h\) 延拓。
\end{solution}

\begin{exercise}\label{b3:ex:16-draft-05}
构造上半平面上两种零边值调和解，并写出保证唯一性的一种附加条件。
\end{exercise}
\begin{solution}
零函数与 \(y\) 都连续取零实轴边值。若要求解有界，则差可奇反射到整平面，反射后的有界调和函数由 Harnack 推出的 Liouville 原理为常数，奇性迫使零。局部反射可由 Poisson 唯一性核验。
\end{solution}

\begin{exercise}\label{b3:ex:16-draft-06}
设有界正则域上 \(\varphi\le\psi\)。证明两个解满足相同的序关系，并估计它们的最大范数差。
\end{exercise}
\begin{solution}
差的边值非负，最大值原理得 \(H_\varphi\le H_\psi\)。把边值差夹在 \(\pm\|\varphi-\psi\|_\infty\) 之间并比较常数解，得内部最大范数差不超过边值差的最大范数。
\end{solution}

\begin{exercise}\label{b3:ex:16-draft-07}
在角宽 \(3\pi/2\) 的扇形顶点构造障碍，说明为何 \(r\cos\theta\) 不适用。
\end{exercise}
\begin{solution}
取 \(\alpha=1/2\)，函数 \(r^{1/2}\cos(\theta/2)\) 在 \(|\theta|\le3\pi/4\) 为正，固定半径上有正下界，到顶点为零。\(r\cos\theta\) 在允许角范围一部分为负，不能作为正障碍。
\end{solution}

\begin{exercise}\label{b3:ex:16-draft-08}
证明有界正则域中任意非空相对开边界片 \(I\) 若承载某个非零非负连续边值，则 \(\omega_D^z(I)>0\) 对所有内部点成立。
\end{exercise}
\begin{solution}
取支集在 \(I\) 的非零非负连续函数 \(\varphi\)。其解非负；若某个内部点为零，则强最小值原理迫使恒零，与边界非零矛盾。因此 \(0<H_\varphi(z)\le\|\varphi\|_\infty\omega_D^z(I)\)。
\end{solution}
